\section{¿Tenemos conexión a Internet?}

\begin{center}
    
  \begin{tcolorbox}[title=Comprobamos la conexión a Internet]
	\orden{\ttfamily ping -c 4 google.es}
  \scriptsize
  \begin{verbatim}
  PING google.es (172.217.17.3) 56(84) bytes de datos.
  64 bytes desde mad07s09-in-f3.1e100.net (172.217.17.3): icmp_seq=1 ttl=119 tiempo=63.4 ms
  64 bytes desde mad07s09-in-f3.1e100.net (172.217.17.3): icmp_seq=2 ttl=119 tiempo=77.9 ms
  64 bytes desde mad07s09-in-f3.1e100.net (172.217.17.3): icmp_seq=3 ttl=119 tiempo=34.3 ms
  64 bytes desde mad07s09-in-f3.1e100.net (172.217.17.3): icmp_seq=4 ttl=119 tiempo=33.3 ms

  --- google.es estadísticas ping ---
  4 paquetes transmitidos, 4 recibidos, 0\% packet loss, time 3004ms
  rtt min/avg/max/mdev = 33.302/52.228/77.929/19.152 ms
  \end{verbatim}
	\end{tcolorbox}

	\begin{tcolorbox}[title=Nota]
		En caso de que sea un portátil o un ordenador con conexión inalámbrica, se deberá usar la aplicación {\color{blue}\ttfamily iwctl}.\\

  \textbf{Pasos a seguir}\\
  Para conectarse a la red WiFi, debemos saber el nombre de nuestro dispositivo primero, para esto un ip link y ya estamos del otro lado. Para los comandos a continuación usaremos \tazul{wlan0} como nombre de dispositivo y supongamos además que nuestra SSID es \tmarron{redCasa} con clave \trojo{1234567890}.

\begin{center}
  \verb+iwctl --passphrase 1234567890 station wlan1 connect redCasa+
\end{center}

Si todo va bien debería funcionar la orden \tmarron{ping}
	\end{tcolorbox}
\end{center}
\newpage
